% ============================================================================
%  CAPÍTULO IV: DESARROLLO DEL TRABAJO
% ============================================================================

\section{Desarrollo del trabajo}

Este capítulo documenta el proceso de configuración del servidor de correo
electrónico sobre Debian, desde la instalación de los paquetes necesarios
hasta las pruebas de conectividad y envío/recepción entre clientes.

% -------------------------------------------------------
\subsection{Arquitectura del sistema}

El servidor de correo se implementó utilizando dos servicios principales
instalados directamente en Debian:

\begin{itemize}
  \item \textbf{Postfix} (MTA): gestiona el envío y recepción de correo
        electrónico vía el protocolo SMTP, habilitando los puertos 25
        (transferencia entre servidores) y 587 (submission para clientes).
  \item \textbf{Dovecot} (MDA): proporciona acceso a los buzones de correo
        mediante los protocolos IMAP (puerto 143) y POP3 (puerto 110),
        además de validar las credenciales de los usuarios a través del
        subsistema SASL.
\end{itemize}

La autenticación de usuarios se realiza a través de Dovecot, que conecta con
Postfix mediante un socket Unix (\texttt{/var/run/dovecot/private/auth}). De
esta forma, el cliente envía sus credenciales al servidor de submission
(puerto 587), Postfix las delega a Dovecot para su validación, y solo tras
la autenticación exitosa se acepta el envío del mensaje.

La figura \ref{fig:arquitectura-correo} presenta la arquitectura general
del servicio de correo implementado.

\begin{figure}[H]
\centering
\begin{tikzpicture}[
    node distance=1.2cm and 1.8cm,
    servicio/.style={draw, rounded corners, minimum width=2.6cm,
                     minimum height=0.9cm, align=center, font=\small},
    cliente/.style={draw, circle, minimum size=1cm, font=\small},
    linea/.style={-Stealth, thick}
  ]
  % Nodos
  \node[servicio] (postfix) {Postfix\\{\footnotesize SMTP}};
  \node[servicio, right=of postfix] (dovecot) {Dovecot\\{\footnotesize IMAP/POP3}};
  \node[cliente, left=of postfix] (c1) {C1};
  \node[cliente, below=of c1] (c2) {C2};
  \node[servicio, below=of postfix, yshift=-0.4cm] (dns) {BIND9\\{\footnotesize DNS}};

  % Conexiones
  \draw[linea] (c1.east) -- node[above,font=\scriptsize]{587} (postfix.west);
  \draw[linea] (c2.east) -| (postfix.west);
  \draw[linea] (postfix.east) -- node[above,font=\scriptsize]{SASL} (dovecot.west);
  \draw[linea, dashed] (postfix.south) -- (dns.north);
\end{tikzpicture}
\caption{Arquitectura del servicio de correo: Postfix, Dovecot y DNS}
\label{fig:arquitectura-correo}
\end{figure}

% -------------------------------------------------------
\subsection{Instalación de paquetes}

Todos los paquetes necesarios se instalaron con el administrador APT de Debian.
El siguiente comando instala las dependencias principales para el servicio de
correo:

\begin{lstlisting}[style=terminal]
$ apt update
$ apt install -y postfix dovecot-core dovecot-lmtpd \
  libsasl2-modules openssl
\end{lstlisting}

\begin{itemize}
  \item \textbf{postfix}: servidor MTA que gestiona el envío y recepción de
        correo vía SMTP.
  \item \textbf{dovecot-core} y \textbf{dovecot-lmtpd}: servidor IMAP/POP3
        con soporte LMTP para la entrega local de correo.
  \item \textbf{openssl}: necesario para la generación de certificados SSL/TLS
        auto-firmados.
  \item \textbf{libsasl2-modules}: módulos SASL para la autenticación de
        usuarios.
\end{itemize}

% -------------------------------------------------------
\subsection{Configuración del servidor DNS (BIND9)}

Antes de configurar el servidor de correo, es necesario contar con un servidor
DNS funcional que resuelva el dominio del proyecto. Se instaló BIND9 con las
siguientes características en la zona directa para \texttt{sudoers.lan}:

\begin{itemize}
  \item Registro MX que apunta al servidor de correo: \texttt{@ IN MX 10 mail.sudoers.lan}.
  \item Registro A para el servidor de correo: \texttt{mail IN A 192.168.0.105}.
  \item Soporte exclusivo IPv4 (\texttt{-4} en el comando \texttt{named}) para
        evitar errores de conectividad con servidores raíz IPv6.
\end{itemize}

La verificación del registro MX se realizó con el comando \texttt{dig}:

\begin{lstlisting}[style=terminal]
$ dig +short MX sudoers.lan
10 mail.sudoers.lan.

$ dig +short mail.sudoers.lan
192.168.0.105
\end{lstlisting}

\capturaAPA{fig:dig-mx}{Verificaci\'on del registro MX del dominio \texttt{sudoers.lan} con \texttt{dig}}{captura-dig-mx.png}{Captura de pantalla mostrando la salida de \texttt{dig +short MX sudoers.lan} y \texttt{dig +short mail.sudoers.lan} en la terminal.}

% -------------------------------------------------------
\subsection{Configuración del servidor SMTP (Postfix)}

La configuración de Postfix se realizó editando el archivo
\texttt{/etc/postfix/main.cf}. A continuación se describen las configuraciones
más relevantes.

\subsubsection{Identidad del servidor}

Se estableció el nombre completo del dominio del correo y el dominio base del
proyecto:

\begin{lstlisting}[style=terminal]
myhostname = mail.sudoers.lan
mydomain   = sudoers.lan
myorigin   = $mydomain
\end{lstlisting}

\subsubsection{Redes permitidas}

El servidor acepta conexiones desde todas las redes privadas, lo que permite
el acceso desde cualquier máquina de la red interna:

\begin{lstlisting}[style=terminal]
mynetworks = 127.0.0.0/8, 10.0.0.0/8,
             172.16.0.0/12, 192.168.0.0/16
\end{lstlisting}

\subsubsection{Almacenamiento de correo}

Se configuró el formato Maildir para el almacenamiento de mensajes, que
almacena cada correo como un archivo individual dentro del directorio
\texttt{Maildir/} del home de cada usuario:

\begin{lstlisting}[style=terminal]
home_mailbox = Maildir/
\end{lstlisting}

\subsubsection{Autenticación SASL}

La autenticación de usuarios se delega a Dovecot mediante el socket SASL.
Esta configuración permite que Postfix valide las credenciales sin gestionar
directamente la base de usuarios:

\begin{lstlisting}[style=terminal]
smtpd_sasl_type           = dovecot
smtpd_sasl_path           = private/auth
smtpd_sasl_auth_enable    = yes
smtpd_sasl_security_options = noanonymous
\end{lstlisting}

\subsubsection{Restricciones de envío}

El servidor acepta envíos únicamente de usuarios autenticados, rechazando
cualquier intento no autorizado:

\begin{lstlisting}[style=terminal]
smtpd_recipient_restrictions =
    permit_sasl_authenticated,
    permit_mynetworks,
    reject_unauth_destination
\end{lstlisting}

\subsubsection{Habilitación de STARTTLS}

Para proteger las conexiones se habilitó STARTTLS con certificados
auto-firmados. El certificado se genera con \texttt{openssl} directamente
en el servidor:

\begin{lstlisting}[style=terminal]
# Generar certificado auto-firmado
$ openssl req -new -x509 -days 3650 -nodes \
    -out /etc/ssl/certs/mail.pem \
    -keyout /etc/ssl/private/mail.key \
    -subj "/CN=mail.sudoers.lan/O=TSO-II/C=AR"
$ chmod 600 /etc/ssl/private/mail.key
\end{lstlisting}

Y en \texttt{/etc/postfix/main.cf} se habilita TLS:

\begin{lstlisting}[style=terminal]
smtpd_use_tls            = yes
smtpd_tls_cert_file      = /etc/ssl/certs/mail.pem
smtpd_tls_key_file       = /etc/ssl/private/mail.key
smtpd_tls_security_level = may
\end{lstlisting}

\subsubsection{Puerto de submission (587)}

El puerto 587 se habilitó en \texttt{/etc/postfix/master.cf} para clientes
que desean enviar correo con autenticación. Esta configuración exige STARTTLS
antes de aceptar autenticación (\texttt{smtpd\_tls\_security\_level=encrypt}):

\begin{lstlisting}[style=terminal]
submission inet n  -  n  -  -  smtpd
  -o smtpd_tls_security_level=encrypt
  -o smtpd_sasl_auth_enable=yes
  -o smtpd_sasl_type=dovecot
  -o smtpd_sasl_path=private/auth
  -o smtpd_relay_restrictions=permit_sasl_authenticated,reject
\end{lstlisting}

La configuración aplicada por Postfix se puede verificar con
\texttt{postconf -n}, que muestra los parámetros efectivos del servidor:

\capturaAPA{fig:postconf}{Par\'ametros efectivos de Postfix con \texttt{postconf -n}}{captura-postconf.png}{Captura de pantalla de la terminal ejecutando \texttt{postconf -n}, mostrando la configuraci\'on activa (myhostname, mynetworks, smtpd\_sasl\_*).}

% -------------------------------------------------------
\subsection{Configuración del servidor IMAP/POP3 (Dovecot)}

Dovecot se configura mediante varios archivos en \texttt{/etc/dovecot/conf.d/}.
Los más relevantes son:

\subsubsection{Protocolos habilitados}

Se habilitaron los protocolos IMAP y POP3 en el archivo
\texttt{10-master.conf}:

\begin{lstlisting}[style=terminal]
protocols = imap pop3
\end{lstlisting}

\subsubsection{Autenticación}

En \texttt{10-auth.conf} se habilitaron los mecanismos de autenticación
\texttt{plain} y \texttt{login}, que permiten la autenticación con
credenciales en texto plano protegidas por STARTTLS:

\begin{lstlisting}[style=terminal]
auth_mechanisms = plain login
\end{lstlisting}

\subsubsection{Formato de buzón}

En \texttt{10-mail.conf} se configuró el formato Maildir, compatible con
la configuración de Postfix:

\begin{lstlisting}[style=terminal]
mail_location = maildir:~/Maildir
\end{lstlisting}

\subsubsection{Certificado SSL}

En \texttt{10-ssl.conf} se habilitó SSL con los certificados auto-firmados
generados previamente:

\begin{lstlisting}[style=terminal]
ssl = required
ssl_cert  = /etc/ssl/certs/mail.pem
ssl_key   = /etc/ssl/private/mail.key
\end{lstlisting}

\capturaAPA{fig:doveconf}{Verificaci\'on de la configuraci\'on de Dovecot con \texttt{doveconf -n}}{captura-doveconf.png}{Captura de pantalla de la terminal ejecutando \texttt{doveconf -n}, mostrando los par\'ametros activos de Dovecot.}

% -------------------------------------------------------
\subsection{Creación de usuarios y estructura de buzón}

Los usuarios de correo se crean como usuarios del sistema en Debian. Cada
usuario necesita una estructura Maildir para almacenar sus mensajes:

\begin{lstlisting}[style=terminal]
# Crear usuarios del sistema
$ useradd -m -s /bin/bash joaquin
$ echo "joaquin:adminpass123" | chpasswd
$ useradd -m -s /bin/bash nicolas
$ echo "nicolas:adminpass123" | chpasswd
$ useradd -m -s /bin/bash david
$ echo "david:adminpass123" | chpasswd

# Crear estructura Maildir para cada usuario
$ for USER in joaquin nicolas david; do
    mkdir -p /home/$USER/Maildir/{cur,new,tmp}
    chown -R $USER:$USER /home/$USER/Maildir
    chmod -R 700 /home/$USER/Maildir
  done
\end{lstlisting}

Una vez creados los usuarios, se inician los servicios de Postfix y Dovecot:

\begin{lstlisting}[style=terminal]
$ systemctl restart postfix
$ systemctl restart dovecot
$ systemctl enable postfix dovecot
\end{lstlisting}

% -------------------------------------------------------
\subsection{Pruebas de conectividad}

Se realizaron pruebas de conectividad desde la máquina cliente hacia el
servidor de correo en \texttt{192.168.0.105}, verificando el estado de
todos los puertos relevantes:

\begin{lstlisting}[style=terminal]
$ ping -c 2 192.168.0.105
2 packets transmitted, 2 received, 0% packet loss

$ for p in 25 110 143 587 993 995; do
    timeout 3 bash -c "echo > /dev/tcp/192.168.0.105/$p" \
      && echo "puerto $p: ABIERTO"
  done
puerto 25: ABIERTO
puerto 110: ABIERTO
puerto 143: ABIERTO
puerto 587: ABIERTO
puerto 993: ABIERTO
puerto 995: ABIERTO
\end{lstlisting}

Todos los puertos del servicio de correo respondieron correctamente,
confirmando que el firewall permite el tráfico y que los servicios
están operativos.

\capturaAPA{fig:puertos-abiertos}{Comprobaci\'on de conectividad a los puertos del servidor de correo}{captura-puertos.png}{Captura de pantalla mostrando la salida del bucle que verifica los puertos 25, 110, 143, 587, 993 y 995 desde la m\'aquina cliente.}

% -------------------------------------------------------
\subsection{Pruebas de envío y recepción de correo}

La verificación del envío de correo se realizó utilizando el protocolo SMTP
directamente, probando la autenticación y el envío entre usuarios del dominio.

\subsubsection{Prueba SMTP básica}

Se verificó que el servidor SMTP responda correctamente con el banner
esperado:

\begin{lstlisting}[style=terminal]
$ echo "QUIT" | nc 192.168.0.105 25
220 mail.sudoers.lan ESMTP Postfix
221 2.0.0 Bye
\end{lstlisting}

\capturaAPA{fig:smtp-banner}{Banner de bienvenida del servidor SMTP (Postfix)}{captura-smtp-banner.png}{Captura de pantalla de la conexi\'on SMTP al puerto 25 mostrando el banner \texttt{220 mail.sudoers.lan ESMTP Postfix}.}

\subsubsection{Prueba IMAP}

Se verificó que Dovecot responda y muestre las capacidades disponibles:

\begin{lstlisting}[style=terminal]
$ echo "LOGOUT" | nc 192.168.0.105 143
* OK [CAPABILITY IMAP4rev1 SASL-IR LOGIN-REFERRALS ID ENABLE
  IDLE LITERAL+ STARTTLS AUTH=PLAIN AUTH=LOGIN]
  Dovecot (Debian) ready.
\end{lstlisting}

\capturaAPA{fig:imap-capabilities}{Capacidades del servidor IMAP (Dovecot)}{captura-imap.png}{Captura de pantalla de la conexi\'on IMAP al puerto 143 mostrando el banner de capacidades de Dovecot.}

\subsubsection{Verificación de STARTTLS}

Se comprobó que el servidor acepte la conexión TLS con certificado
auto-firmado:

\begin{lstlisting}[style=terminal]
$ openssl s_client -connect 192.168.0.105:587 -starttls smtp
subject=CN=mail.sudoers.lan, O=TSO-II, C=AR
Verify return code: 18 (self-signed certificate)
\end{lstlisting}

\capturaAPA{fig:starttls}{Negociaci\'on STARTTLS con el servidor SMTP}{captura-starttls.png}{Captura de pantalla de \texttt{openssl s\_client} mostrando el sujeto del certificado y el c\'odigo de verificaci\'on 18 (self-signed).}

\subsubsection{Envío con autenticación (swaks)}

Para una prueba completa de envío se utilizó la herramienta \texttt{swaks}
(\textit{Swiss Army Knife for SMTP}):

\begin{lstlisting}[style=terminal]
$ swaks --to joaquin@sudoers.lan \
    --from nicolas@sudoers.lan \
    --server 192.168.0.105 --port 587 \
    --auth LOGIN --auth-user nicolas \
    --auth-password adminpass123 \
    --header "Subject: Prueba" \
    --body "Hola Joaquin, esto es una prueba."
\end{lstlisting}

La figura \ref{fig:swaks} muestra el resultado exitoso del envío.

\capturaAPA{fig:swaks}{Env\'io de correo autenticado con \texttt{swaks}}{captura-swaks.png}{Captura de pantalla de \texttt{swaks} mostrando el intercambio SMTP, la autenticaci\'on exitosa y el c\'odigo 250 para la aceptaci\'on del mensaje.}

La respuesta del servidor confirmó el envío exitoso del mensaje. Para
verificar la recepción, se puede inspeccionar el buzón del destinatario:

\begin{lstlisting}[style=terminal]
$ ls /home/joaquin/Maildir/new/
1627...joaquin@sudoers.lan
\end{lstlisting}

\capturaAPA{fig:maildir-receptor}{Buz\'on Maildir del destinatario con el mensaje recibido}{captura-maildir.png}{Captura de pantalla listando el contenido de \texttt{/home/joaquin/Maildir/new/} donde aparece el correo recibido.}

% -------------------------------------------------------
\subsection{Acceso desde clientes}

Para acceder al correo desde un cliente como Thunderbird, se utilizó la
siguiente configuración:

\begin{table}[H]
\centering
\begin{tabular}{ll}
\toprule
\textbf{Parámetro} & \textbf{Valor} \\
\midrule
Servidor de entrada (IMAP) & mail.sudoers.lan \\
Puerto de entrada & 143 \\
Seguridad de entrada & STARTTLS \\
Servidor de salida (SMTP) & mail.sudoers.lan \\
Puerto de salida & 587 \\
Seguridad de salida & STARTTLS \\
Autenticación & Contraseña normal \\
Usuario & joaquin \\
\bottomrule
\end{tabular}
\caption{Configuración del cliente de correo Thunderbird}
\label{tab:config-thunderbird}
\end{table}

Debido a que el certificado es auto-firmado, el cliente mostrará un aviso
de seguridad la primera vez que se conecte. Es necesario aceptar la excepción
de seguridad manualmente para permitir la conexión.

\capturaAPA{fig:thunderbird}{Cliente de correo Thunderbird conectado al servidor IMAP con los mensajes recibidos}{captura-thunderbird.png}{Captura de pantalla del cliente Thunderbird mostrando la cuenta joaquin@sudoers.lan conectada y la bandeja de entrada con el mensaje de prueba recibido.}

% -------------------------------------------------------
\subsection{Usuarios creados}

El servidor de correo cuenta con los siguientes usuarios, creados como
usuarios del sistema Debian:

\begin{table}[H]
\centering
\begin{tabular}{lll}
\toprule
\textbf{Usuario} & \textbf{Contraseña} & \textbf{Correo} \\
\midrule
joaquin & adminpass123 & joaquin@sudoers.lan \\
nicolas & adminpass123 & nicolas@sudoers.lan \\
david & adminpass123 & david@sudoers.lan \\
\bottomrule
\end{tabular}
\caption{Usuarios del servidor de correo}
\label{tab:usuarios-correo}
\end{table}

\begin{quote}
\small\textit{Nota:} las contraseñas indicadas son las utilizadas en el
entorno de laboratorio. En un entorno de producción se recomienda utilizar
contraseñas robustas y rotarlas periódicamente.
\end{quote}
